% 辛群（综述）
% license CCBY4
% type Wiki

本文根据 CC-BY-SA 协议转载翻译自维基百科\href{https://en.wikipedia.org/wiki/Symplectic_group}{相关文章}

在数学中，“辛群”的名称可以指两类不同但密切相关的数学群族，分别记作Sp(2n, F)与Sp(n)，其中n正整数，F为域（通常取C或 R）。后一种被称为紧致辛群，也记作$\mathrm{USp}(n)$。许多作者使用略有不同的记号，通常差别体现在系数 2 的处理上。此处采用的记号与表示这些群的最常见矩阵的规模保持一致。在卡当（Cartan）对单李代数的分类中，复群Sp(2n, C) 的李代数记作Cₙ，而Sp(n)则是Sp(2n, C)的紧致实型。需要注意的是，当我们提及“（紧致）辛群”时，通常是指依照维数n编号的一族（紧致）辛群。

“辛群”这一名称由赫尔曼·外尔提出，用来取代此前令人困惑的名称——“（线）复群”（(line) complex group）和“阿贝尔线性群”。它是“复数”一词的希腊语类比。

梅塔辛群是实数域上辛群的双覆盖；它在其他局部域、有限域以及阿德尔环上也有类似的对应物。
\subsection{Sp(2n, F)}
辛群是一个经典群，定义为作用在2n 维向量空间（定义在域$F$上）上的线性变换的集合，这些变换保持一个非退化的斜对称双线性型。这样的向量空间称为辛向量空间，而一个抽象辛向量空间$V$的辛群记作$\mathrm{Sp}(V)$。一旦为$V$选定基底，辛群就变为由矩阵乘法构成的$2n \times 2n$辛矩阵群，其元素属于域$F$。这个群记作$\mathrm{Sp}(2n, F)$或$\mathrm{Sp}(n, F)$。如果该双线性型由一个非奇异斜对称矩阵 $\Omega$ 表示，则有：
$$
\operatorname{Sp}(2n,F) = \{ M \in M_{2n \times 2n}(F) : M^{\mathrm{T}} \, \Omega \, M = \Omega \},~
$$
其中$M^{\mathrm{T}}$表示矩阵$M$的转置。通常，$\Omega$定义为：
$$
\Omega = 
\begin{pmatrix}
0 & I_n \\
- I_n & 0
\end{pmatrix},~
$$
其中$I_n$是 $n \times n$ 单位矩阵。在这种情况下，$\mathrm{Sp}(2n, F)$ 可以表示为如下的分块矩阵：
$$
\begin{pmatrix}
A & B \\
C & D
\end{pmatrix}, \quad A, B, C, D \in M_{n \times n}(F),~
$$
并满足以下三个方程：
$$
\begin{aligned}
- C^{\mathrm{T}} A + A^{\mathrm{T}} C &= 0, \\
- C^{\mathrm{T}} B + A^{\mathrm{T}} D &= I_n, \\
- D^{\mathrm{T}} B + B^{\mathrm{T}} D &= 0.
\end{aligned}~
$$
由于所有辛矩阵的行列式都等于 1，辛群是特殊线性群$\mathrm{SL}(2n, F)$的一个子群。当$n = 1$时，矩阵满足辛条件当且仅当其行列式为 1，因此：$\mathrm{Sp}(2, F) = \mathrm{SL}(2, F)$.而当$n > 1$时，还存在附加条件，即$\mathrm{Sp}(2n, F)$是$\mathrm{SL}(2n, F)$的一个真子群。

通常，域$F$取实数域$\mathbf{R}$或复数域$\mathbf{C}$。在这些情况下，$\mathrm{Sp}(2n, F)$分别是实或复的李群，其实或复维数为：$n(2n+1)$.这些群是连通的，但不是紧致的。

当域的特征不为 2 时，$\mathrm{Sp}(2n, F)$的中心由矩阵$I_{2n}$ 和 $-I_{2n}$ 构成。\(^\text{[1]}\)由于其中心是离散的，且模去中心的商群是单群，因此 $\mathrm{Sp}(2n, F)$ 被视为一个单李群。

对应李代数的实秩，因此李群$\mathrm{Sp}(2n, F)$的实秩为：$n$。

辛群$\mathrm{Sp}(2n, F)$的李代数定义为：
$$
\mathfrak{sp}(2n,F) = \{ X \in M_{2n \times 2n}(F) : \Omega X + X^{\mathrm{T}} \Omega = 0 \},~
$$
其李括号由对易子给出。\(^\text{[2]}\)

对于标准的斜对称双线性型：$\Omega = \begin{pmatrix}0 & I \\- I & 0
\end{pmatrix}$，该李代数由所有分块矩阵$\begin{pmatrix}A & B \\C & D
\end{pmatrix}$构成，并满足条件：
$$
\begin{aligned}
A &= - D^{\mathrm{T}}, \\
B &= B^{\mathrm{T}}, \\
C &= C^{\mathrm{T}}
\end{aligned}~
$$
\subsubsection{Sp(2n, C)}
复数域上的辛群是一个非紧致、单连通的单李群。该群的定义中并不涉及共轭（与直觉上的预期相反），而只是将定义中的域换为复数域，其余保持完全相同。\(^\text{[3]}\)
\subsubsection{Sp(2n, R)}
$\mathrm{Sp}(n, \mathbf{C})$是实群$\mathrm{Sp}(2n, \mathbf{R})$ 的复化。$\mathrm{Sp}(2n, \mathbf{R})$是一个实的、非紧致的、连通的单李群。\(^\text{[4]}\)它的基本群同构于加法下的整数群。作为单李群的实型，它的李代数是可分解李代数。

$\mathrm{Sp}(2n, \mathbf{R})$的进一步性质：
\begin{itemize}
\item 从李代数$\mathfrak{sp}(2n, \mathbf{R})$ 到群 $\mathrm{Sp}(2n, \mathbb{R})$的指数映射不是满射。然而，该群的任意元素都可以表示为两个指数的乘积。\(^\text{[5]}\)换句话说：
   $$
   \forall S \in \operatorname{Sp}(2n,\mathbf{R}) \;\; \exists X, Y \in \mathfrak{sp}(2n,\mathbf{R}) \;\; S = e^{X} e^{Y}.~
   $$
\item 对于所有$S \in \mathrm{Sp}(2n, \mathbf{R})$，有：
   $$
   S = O Z O',~
   $$
   其中
   $$
   O, O' \in \operatorname{Sp}(2n, \mathbf{R}) \cap SO(2n) \cong U(n),~
   $$
   且
   $$
   Z = \begin{pmatrix} D & 0 \\ 0 & D^{-1} \end{pmatrix}.~
   $$
   这里 $D$ 是正定对角矩阵。所有这样的 $Z$ 构成 $\mathrm{Sp}(2n, \mathbf{R})$ 的一个非紧子群，而 $U(n)$ 构成一个紧子群。这种分解被称为欧拉分解或Bloch–Messiah 分解。\(^\text{[6]}\)更多关于辛矩阵的性质可参见维基百科相关页面。
\item 作为李群，$\mathrm{Sp}(2n, \mathbf{R})$具有流形结构。该流形与酉群$U(n)$和一个维数为$n(n+1)$的向量空间的笛卡尔积微分同胚。\(^\text{[7]}\)
\end{itemize}
\subsubsection{无穷小生成元}
辛李代数$\mathfrak{sp}(2n, F)$的元素称为哈密顿矩阵。

这些矩阵$Q$的一般形式为：
$$
Q = 
\begin{pmatrix}
A & B \\
C & -A^{\mathrm{T}}
\end{pmatrix},~
$$
其中$B$与$C$是对称矩阵。推导过程可参见“经典群”。
\subsubsection{辛矩阵的例子}
对于$\mathrm{Sp}(2, \mathbf{R})$，即所有行列式为 1 的$2 \times 2$ 矩阵，该群的三个辛$(0,1)$-矩阵为：\(^\text{[8]}\)
$$
\begin{pmatrix}
1 & 0 \\
0 & 1
\end{pmatrix}, \quad
\begin{pmatrix}
1 & 0 \\
1 & 1
\end{pmatrix}, \quad
\text{以及} \quad
\begin{pmatrix}
1 & 1 \\
0 & 1
\end{pmatrix}.~
$$
\textbf{Sp(2n, R)}

事实证明，$\operatorname{Sp}(2n, \mathbf{R})$可以用生成元作出相当明确的描述。设$\operatorname{Sym}(n)$表示所有对称的$n \times n$矩阵，则$\operatorname{Sp}(2n, \mathbf{R})$由以下集合生成：$D(n) \cup N(n) \cup \{\Omega\}$,其中
$$
\begin{aligned}
D(n) &= \left\{ 
\begin{bmatrix}
A & 0 \\
0 & (A^{T})^{-1}
\end{bmatrix}
\;\middle|\; A \in \operatorname{GL}(n, \mathbf{R})
\right\},\\
N(n) &= \left\{
\begin{bmatrix}
I_{n} & B \\
0 & I_{n}
\end{bmatrix}
\;\middle|\; B \in \operatorname{Sym}(n)
\right\},
\end{aligned}~
$$
它们都是 $\operatorname{Sp}(2n, \mathbf{R})$ 的子群。\(^\text{[9] p.173 [10] p.2}\)
\subsubsection{与辛几何的关系}
辛几何研究的是辛流形。在辛流形上的任意一点，其切空间都是一个辛向量空间。\(^\text{[11]}\)如前所述，保持辛向量空间结构的变换构成一个群，这个群就是$\mathrm{Sp}(2n, F)$，具体形式依赖于空间的维数以及其定义的域。

一个辛向量空间本身就是一个辛流形。因此，在辛群作用下的变换，在某种意义上可以看作是辛同胚的线性化版本，而辛同胚是辛流形上一类更一般的结构保持变换。
\subsection{Sp(n)}
紧致辛群\(^\text{[12]}\)$\mathrm{Sp}(n)$ 定义为复辛群$\mathrm{Sp}(2n, \mathbb{C})$与$2n \times 2n$酉群的交集：
$$
\operatorname{Sp}(n) := \operatorname{Sp}(2n;\mathbb{C}) \cap U(2n) 
= \operatorname{Sp}(2n;\mathbb{C}) \cap SU(2n).~
$$
它有时记作$\mathrm{USp}(2n)$。

另一种等价的描述是：$\mathrm{Sp}(n)$可以看作是$\mathrm{GL}(n, \mathbb{H})$（可逆四元数矩阵群）的一个子群，它保持$\mathbb{H}^n$上的标准 Hermitian 型：
$$
\langle x, y \rangle = \bar{x}_1 y_1 + \cdots + \bar{x}_n y_n~
$$
换句话说，$\mathrm{Sp}(n)$就是四元数酉群$U(n, \mathbb{H})$。\(^\text{[13]}\)它有时也被称为超酉群。此外，$\mathrm{Sp}(1)$即为所有模长为 1 的四元数的集合，与$\mathrm{SU}(2)$等价，并且在拓扑上同胚于三维球面$S^3$。

需要注意的是，$\mathrm{Sp}(n)$并不是前文所述意义上的“辛群”——它并不保持 $\mathbb{H}^n$上的某个非退化斜对称$\mathbb{H}$-双线性型，因为除了零型之外不存在这样的形式。相反，它同构于$\mathrm{Sp}(2n, \mathbb{C})$ 的一个子群，因此确实保持一个在复向量空间（维数为$2n$的两倍）上的复辛形式。正如下文所解释的，$\mathrm{Sp}(n)$ 的李代数是复辛李代数 $\mathfrak{sp}(2n, \mathbb{C})$ 的紧实型。

$\mathrm{Sp}(n)$ 是一个实李群，其（实）维数为：$n(2n+1)$。它是紧致的并且是单连通的。[14]

$\mathrm{Sp}(n)$的李代数由四元数斜 Hermitian 矩阵给出，即所有满足条件的$n \times n$四元数矩阵：
$$
A + A^{\dagger} = 0,~
$$
其中$A^{\dagger}$表示$A$的共轭转置（此处取四元数共轭）。其李括号由对易子给出。
\subsubsection{重要子群}
一些主要的子群关系为：
$$
\begin{aligned}
\operatorname{Sp}(n) &\supset \operatorname{Sp}(n-1)\\
\operatorname{Sp}(n) &\supset \operatorname{U}(n)\\
\operatorname{Sp}(2) &\supset \operatorname{O}(4)
\end{aligned}~
$$
反过来，$\mathrm{Sp}(n)$ 本身也是其他一些群的子群：
$$
\begin{aligned}
\operatorname{SU}(2n) &\supset \operatorname{Sp}(n)\\
F_{4} &\supset \operatorname{Sp}(4)\\
G_{2} &\supset \operatorname{Sp}(1)
\end{aligned}~
$$
此外，还有以下李代数之间的同构关系：$\mathfrak{sp}(2) \cong \mathfrak{so}(5), \quad\mathfrak{sp}(1) \cong \mathfrak{so}(3) \cong \mathfrak{su}(2)$。
\subsection{辛群之间的关系}
每个复半单李代数都有一个分裂实型和一个紧实型；前者可以看作是后两者的复化。

$\mathrm{Sp}(2n, \mathbf{C})$的李代数是半单的，记作$\mathfrak{sp}(2n, \mathbf{C})$。它的分裂实型是$\mathfrak{sp}(2n, \mathbf{R})$，而它的紧实型是$\mathfrak{sp}(n)$。这两者分别对应于李群$\mathrm{Sp}(2n, \mathbf{R})$和$\mathrm{Sp}(n)$。

李代数$\mathfrak{sp}(p, n-p)$（对应于李群$\mathrm{Sp}(p, n-p)$）则是紧实型的不定号对应物。
\subsection{物理意义}
\subsubsection{经典力学}
非紧致辛群$\mathrm{Sp}(2n, \mathbb{R})$在经典物理中出现，作为保持泊松括号的正则坐标对称性。

考虑一个含有$n$个粒子的系统，它按照哈密顿方程演化。在给定时刻，它在相空间中的位置由正则坐标向量表示为：
$$
\mathbf{z} = (q^{1}, \ldots, q^{n}, p_{1}, \ldots, p_{n})^{\mathrm{T}}~
$$
群$\mathrm{Sp}(2n, \mathbb{R})$的元素在某种意义上就是该向量上的正则变换，即它们保持哈密顿方程的形式。\(^\text{[15][16]}\)

若新的正则坐标为：
$$
\mathbf{Z} = \mathbf{Z}(\mathbf{z}, t) = (Q^{1}, \ldots, Q^{n}, P_{1}, \ldots, P_{n})^{\mathrm{T}},~
$$
则用“点”表示对时间的导数，有：
$$
\dot{\mathbf{Z}} = M(\mathbf{z}, t) \, \dot{\mathbf{z}},~
$$
其中：
$$
M(\mathbf{z}, t) \in \operatorname{Sp}(2n, \mathbf{R})~
$$
对所有时刻$t$以及相空间中的所有$\mathbf{z}$都成立。\(^\text{[17]}\)

对于黎曼流形这一特殊情形，哈密顿方程描述了该流形上的测地线。坐标$q^{i}$位于底流形上，而动量$p_{i}$则位于余切丛中。这就是为什么习惯上将它们分别写作上标和下标，以区分它们所在的位置。对应的哈密顿量完全由动能组成：$H = \tfrac{1}{2} g^{ij}(q) p_{i} p_{j}$,其中 $g^{ij}$ 是黎曼流形度量张量 $g_{ij}$ 的逆。\(^\text{[18][16]}\)事实上，任何光滑流形的余切丛都可以以一种正则方式赋予辛结构，其中的辛形式定义为典范一形式的外微分。\(^\text{[19]}\)
\subsubsection{量子力学}
考虑一个包含$n$个粒子的系统，其量子态同时编码了位置和动量。这些坐标是连续变量，因此状态所在的希尔伯特空间是无限维的。这往往使得分析该情形变得复杂。另一种方法是在相空间中研究位置和动量算符在海森堡方程下的演化。

构造正则坐标向量：
$$
\mathbf{\hat{z}} = (\hat{q}^{1}, \ldots, \hat{q}^{n}, \hat{p}_{1}, \ldots, \hat{p}_{n})^{\mathrm{T}}~
$$
正则对易关系可以简洁地表示为：
$$
[\mathbf{\hat{z}}, \mathbf{\hat{z}}^{\mathrm{T}}] = i \hbar \, \Omega~
$$
其中：
$$
\Omega = 
\begin{pmatrix}
\mathbf{0} & I_{n} \\
- I_{n} & \mathbf{0}
\end{pmatrix},~
$$
而$I_n$是$n \times n$的单位矩阵。

许多物理情形只需要考虑二次型哈密顿量，即具有如下形式的哈密顿量：
$$
\hat{H} = \tfrac{1}{2} \mathbf{\hat{z}}^{\mathrm{T}} K \mathbf{\hat{z}},~
$$
其中$K$是一个$2n \times 2n$的实对称矩阵。这一限制非常有用，它允许我们将海森堡方程改写为：
$$
\frac{d \mathbf{\hat{z}}}{dt} = \Omega K \mathbf{\hat{z}}~
$$
该方程的解必须保持正则对易关系。可以证明，该系统的时间演化等价于实辛群 $\mathrm{Sp}(2n, \mathbb{R})$在相空间上的作用。
\subsection{参见}
\begin{itemize}
\item 哈密顿力学 
\item 梅塔辛群 
\item 正交群
\item 参数模群
\item 射影酉群 
\item 经典李群的表示 
\item 辛流形、辛矩阵、辛向量空间、辛表示
\item 酉群 
\item Θ10
\end{itemize}
\subsection{注释}
\begin{enumerate}
\item “Symplectic group”，Encyclopedia of Mathematics，检索于2014年12月13日。
\item Hall 2015, 命题 3.25。
\item Hall 2015，第 10 页。
\item “Is the symplectic group Sp(2n, R) simple?”，Stack Exchange，检索于2014年12月14日。
\item “Is the exponential map for Sp(2n, R) surjective?”，Stack Exchange，检索于2014年12月5日。
\item “Standard forms and entanglement engineering of multimode Gaussian states under local operations – Serafini and Adesso”，检索于2015年1月30日。
\item “Symplectic Geometry – Arnol'd and Givental”，检索于2015年1月30日。
\item Symplectic Group，（来源：Wolfram MathWorld），下载于2012年2月14日。
\item Gerald B. Folland. (2016). Harmonic analysis in phase space。Princeton: Princeton Univ Press，第173页。ISBN 978-1-4008-8242-7。OCLC 945482850。
\item Habermann, Katharina, 1966- (2006). Introduction to symplectic Dirac operators。Springer。ISBN 978-3-540-33421-7。OCLC 262692314。
\item “Lecture Notes – Lecture 2: Symplectic reduction”，检索于2015年1月30日。
\item Hall 2015，第 1.2.8 节。
\item Hall 2015，第 14 页。
\item Hall 2015, 命题 13.12。
\item Arnold 1989 对经典力学进行了广泛的数学综述。参见第8章关于辛流形的讨论。
\item Ralph Abraham and Jerrold E. Marsden, Foundations of Mechanics, (1978) Benjamin-Cummings, London。ISBN 0-8053-0102-X。
\item Goldstein 1980，第 9.3 节。
\item Jürgen Jost, (1992) Riemannian Geometry and Geometric Analysis, Springer。
\item da Silva, Ana Cannas (2008). Lectures on Symplectic Geometry. Lecture Notes in Mathematics, Vol. 1764. Berlin, Heidelberg: Springer Berlin Heidelberg，第9页。doi:10.1007/978-3-540-45330-7。ISBN 978-3-540-42195-5。
\end{enumerate}
\subsection{参考文献}
\begin{itemize}
\item Arnold, V. I. (1989). Mathematical Methods of Classical Mechanics Graduate Texts in Mathematics, Vol. 60 (第二版)，Springer-Verlag。ISBN 0-387-96890-3。
\item Hall, Brian C. (2015). Lie groups, Lie algebras, and representations: An elementary introduction, Graduate Texts in Mathematics, Vol. 222 (第二版)，Springer。ISBN 978-3319134666。
\item Fulton, W.; Harris, J. (1991). Representation Theory, A First Course, Graduate Texts in Mathematics, Vol. 129, Springer-Verlag。ISBN 978-0-387-97495-8。
\item Goldstein, H. (1980) [1950]. “第七章”，Classical Mechanics (第二版)，Reading MA: Addison-Wesley。ISBN 0-201-02918-9。
\item Lee, J. M. (2003). Introduction to Smooth Manifolds, Graduate Texts in Mathematics, Vol. 218, Springer-Verlag。ISBN 0-387-95448-1。
\item Rossmann, Wulf (2002). Lie Groups – An Introduction Through Linear Groups, Oxford Graduate Texts in Mathematics, Oxford Science Publications。ISBN 0-19-859683-9。
\item Ferraro, Alessandro; Olivares, Stefano; Paris, Matteo G. A. (2005年3月). “Gaussian states in continuous variable quantum information”，arXiv:quant-ph/0503237。
\end{itemize}
